\section*{Chapter \ref{ch:g8:speed} --- Proportionality, Speed and Averages}

\begin{solution}{exo:g8:speed:1}
$\frac{150}{2} = 75$ km/h. \quad
$45$ min $= 0.75$ h: $\frac{27}{0.75} = 36$ km/h. \quad
$\frac{100}{12.5} = 8$ m/s.
\end{solution}

\begin{solution}{exo:g8:speed:2}
$2 \times 85 = 170$ km. \quad
$40$ min $= \frac23$ h: $90 \times \frac23 = 60$ km. \quad
$\frac{315}{90} = 3.5$ h $= 3$ h $30$ min.
\end{solution}

\begin{solution}{exo:g8:speed:3}
$72 \div 3.6 = 20$ m/s. \quad
$5 \times 3.6 = 18$ km/h. \quad
$108 \div 3.6 = 30$ m/s.
\end{solution}

\begin{solution}{exo:g8:speed:4}
$600 \div 15 = 40$ min. \quad
$2$ h $30 = 150$ min: $15 \times 150 = 2250$ L.
\end{solution}

\begin{solution}{exo:g8:speed:5}
$\frac{3.30}{1.2} = 2.75$ euros/kg against
$\frac{2.16}{0.8} = 2.70$ euros/kg: the second offer is (slightly)
the better buy.
\end{solution}

\begin{solution}{exo:g8:speed:6}
\[
\bar v = \frac{25 \times 12.4 + 35 \times 10}{60}
= \frac{310 + 350}{60} = \frac{660}{60} = 11 .
\]
\end{solution}

\begin{solution}{exo:g8:speed:7}
$d = 340 \times 6 = 2040$ m $\approx 2$ km.
\end{solution}

\begin{solution}{exo:g8:speed:8}
Distances: $2 \times 5 = 10$ km, then $1.5 \times 4 = 6$ km: total
$16$ km. Elapsed time: $2 + 0.5 + 1.5 = 4$ h. \omterm{def:g8:speed:def}{Average speed}:
$\frac{16}{4} = 4$ km/h.
\end{solution}

\begin{solution}{exo:g8:speed:9}
Weighted:
$\frac{3 \times 15 + 2 \times 9 + 1 \times 12}{3 + 2 + 1}
= \frac{45 + 18 + 12}{6} = \frac{75}{6} = 12.5$.
Unweighted: $\frac{15 + 9 + 12}{3} = 12$. They differ because the
coefficients give the mark $15$ three times the weight of the mark
$12$.
\end{solution}

\begin{solution}{exo:g8:speed:10}
\emph{1.} Times: $\frac{120}{80} = 1.5$ h and $\frac{120}{120} = 1$
h. \omterm{def:g8:speed:def}{Average speed}: $\frac{240}{2.5} = 96$ km/h.

\emph{2.} The car spends \emph{more time} at $80$ km/h ($1.5$ h)
than at $120$ km/h ($1$ h): the slow \omterm{def:g8:speed:def}{speed} weighs more in the
time-weighted average, pulling it below the midpoint $100$.
\end{solution}

\begin{solution}{exo:g8:speed:11}
The gap of $18$ km closes at $5 + 4 = 9$ km/h: they meet after
$\frac{18}{9} = 2$ h, at $12{:}00$. Anna has then walked
$2 \times 5 = 10$ km: they meet $10$ km from Anna's village (and
$8$ km from Boris's --- check: $10 + 8 = 18$).
\end{solution}

\begin{solution}{pb:g8:speed:1}
\textbf{1.} By \cref{def:g8:speed:def}, $t_1 = \frac{d}{v_1}$ and
$t_2 = \frac{d}{v_2}$; the whole journey lasts
$\frac{d}{v_1} + \frac{d}{v_2}$ for a distance of $2d$.

\textbf{2.} Factor $d$ out of the \omterm{def:g4:fractions:def}{denominator} and simplify it
away:
\[
\bar v
= \frac{2d}{d\left(\frac{1}{v_1} + \frac{1}{v_2}\right)}
= \frac{2}{\ \frac{1}{v_1} + \frac{1}{v_2}\ } .
\]
Common \omterm{def:g4:fractions:def}{denominator} inside:
$\frac{1}{v_1} + \frac{1}{v_2} = \frac{v_2 + v_1}{v_1 v_2}$, and
dividing means multiplying by the \omterm{def:g8:fractions:inverse}{inverse}
(\cref{thm:g8:fractions:div}):
\[
\bar v = 2 \times \frac{v_1 v_2}{v_1 + v_2}
= \frac{2\, v_1 v_2}{v_1 + v_2} .
\]

\textbf{3.} $\dfrac{2 \times 30 \times 15}{30 + 15}
= \dfrac{900}{45} = 20$ km/h, as in
\cref{ex:g8:speed:twolegs}; and
$\dfrac{2 \times 80 \times 120}{80 + 120} = \dfrac{19\,200}{200}
= 96$ km/h, as found in \cref{exo:g8:speed:10}.

\textbf{4.} The average depends only on the two \omterm{def:g8:speed:def}{speeds}, not on
$d$: a $1$ km round trip and a $100$ km round trip at the same
two \omterm{def:g8:speed:def}{speeds} have exactly the same \omterm{def:g8:speed:def}{average speed}. (That is why the
formula could be checked on journeys of different lengths.)

\textbf{5.} In time $T$ at each \omterm{def:g8:speed:def}{speed}, the distances are $v_1 T$
and $v_2 T$, so
\[
\bar v = \frac{v_1 T + v_2 T}{2T}
= \frac{(v_1 + v_2)\,T}{2T}
= \frac{v_1 + v_2}{2} .
\]
An average of \omterm{def:g8:speed:def}{speeds} is always weighted by \emph{time}
(\cref{def:g8:speed:average}). With equal times the two \omterm{def:g8:speed:def}{speeds}
carry equal weights: plain midpoint. With equal distances the
slower \omterm{def:g8:speed:def}{speed} occupies \emph{more} time, hence more weight ---
and the average slides toward it.

\textbf{6.} For $(30, 15)$: $H = 20$ and $A = 22.5$. For
$(80, 120)$: $H = 96$ and $A = 100$. The arithmetic mean wins
both times.

\textbf{7.} Over the common \omterm{def:g4:fractions:def}{denominator} $2(v_1 + v_2)$:
\[
A - H
= \frac{(v_1 + v_2)^2 - 4\, v_1 v_2}{2\,(v_1 + v_2)} ,
\]
and the \omterm{def:g4:fractions:def}{numerator} expands with the identities of
\cref{pb:g8:equations:1}:
\[
(v_1 + v_2)^2 - 4 v_1 v_2
= v_1^2 + 2 v_1 v_2 + v_2^2 - 4 v_1 v_2
= v_1^2 - 2 v_1 v_2 + v_2^2
= (v_1 - v_2)^2 .
\]

\textbf{8.} A square is never negative, and $2(v_1 + v_2)$ is
positive: so $A - H \geq 0$, that is $H \leq A$, for every pair
of positive \omterm{def:g8:speed:def}{speeds}. Equality requires $(v_1 - v_2)^2 = 0$, i.e.\
$v_1 = v_2$: the moment the two \omterm{def:g8:speed:def}{speeds} differ, the round trip's
average drops strictly below the midpoint --- the surprise is a
theorem.

\textbf{9.} Multiply:
\[
H \times A
= \frac{2\, v_1 v_2}{v_1 + v_2} \times \frac{v_1 + v_2}{2}
= v_1 v_2 ,
\]
after cancelling the factor $2$ and the factor $v_1 + v_2$.

\textbf{10.} Formula: $\bar v = \dfrac{2 \times 40 \times 60}
{40 + 60} = \dfrac{4800}{100} = 48$ km/h. The long way:
$\frac{60}{40} = 1.5$ h and $\frac{60}{60} = 1$ h, so $120$ km in
$2.5$ h: $\frac{120}{2.5} = 48$ km/h. They agree.

\textbf{11.} Averaging $30$ km/h over $30$ km means a total time
of at most $\frac{30}{30} = 1$ h. The first $15$ km at $15$ km/h
already take $\frac{15}{15} = 1$ h: the \emph{entire} time budget
is spent, with \omterm{def:g3:division:half}{half} the course to go.

\textbf{12.} At $90$ km/h the second \omterm{def:g3:division:half}{half} takes
$\frac{15}{90} = \frac16$ h, so the race takes $\frac76$ h and
\[
\bar v = \frac{30}{\ 7/6\ } = \frac{180}{7} \approx 25.7
\text{ km/h} < 30 .
\]
Whatever the \omterm{def:g3:division:half}{second-half} \omterm{def:g8:speed:def}{speed}, its duration is a positive
number, so the total time exceeds $1$ h and the average stays
strictly below $30$ km/h. (With the formula: $\frac{2 \times 15
\times v_2}{15 + v_2} = 30$ would force $30 v_2 = 450 + 30 v_2$,
i.e.\ $0 = 450$ --- no \omterm{def:g8:speed:def}{speed} $v_2$ works.)

\textbf{13.} A $20$ km/h average allows $\frac{30}{20} = 1.5$ h;
after the first hour, $15$ km remain for $0.5$ h: she needs
$\frac{15}{0.5} = 30$ km/h. Check:
$H(15, 30) = \frac{2 \times 15 \times 30}{45} = \frac{900}{45}
= 20$ km/h.

\textbf{14.} Per minute, the taps fill $\frac{1}{20}$ and
$\frac{1}{30}$ of the tub; together
\[
\frac{1}{20} + \frac{1}{30}
= \frac{3}{60} + \frac{2}{60}
= \frac{5}{60} = \frac{1}{12}
\]
of the tub per minute: the tub fills in $12$ min. And
$H(20, 30) = \frac{2 \times 600}{50} = 24 = 2 \times 12$: the two
taps together take exactly \emph{\omterm{def:g3:division:half}{half} the \omterm{pb:g8:speed:1}{harmonic mean}} of their
solo times --- rates add, so it is again harmonic-mean territory.

\textbf{15.} Same distance each way, so
\[
\bar v = \frac{2 \times 900 \times 700}{900 + 700}
= \frac{1\,260\,000}{1\,600} = 787.5 \text{ km/h},
\]
below the still-air $800$ km/h. In one sentence: the plane spends
\emph{longer} flying against the wind than with it, so the slow
leg gets the larger weight in the time-weighted average --- a
wind can only lengthen a round trip.
\end{solution}
